\documentclass[10pt,a4paper]{amsart}
\usepackage[margin=19mm]{geometry}
\usepackage{amsmath,amssymb,mathtools}
\usepackage[scr=rsfs,cal=euler]{mathalfa}
\usepackage{xfrac}
\usepackage[unicode,hidelinks,bookmarks=false]{hyperref}
\newcommand{\ie}{i.e.\ }
\newcommand{\cf}{cf.\ }
\newcommand{\resp}{respectively\ }
\newcommand{\Subsection}{\subsection}
\providecommand{\nobreakdash}{\nobreak\hbox{-}}
\setlength{\emergencystretch}{2em}
\multlinegap0pt
\usepackage{mathtools}
\usepackage{amssymb}
\usepackage{tikz}
\usepackage[shortlabels]{enumitem}
\usepackage{xcolor}
\newtheorem{lemma}{Lemma}
\newtheorem{theorem}{Theorem}
\usepackage{cleveref}
\crefname{lemma}{Lemma}{Lemmas}
\crefname{theorem}{Theorem}{Theorems}
\newcommand{\D}{\mathrm{D}}
\newcommand{\I}{\mathrm{I}}
\newcommand{\M}{\mathbb{M}}
\newcommand{\Mcal}{\mathcal{M}}
\newcommand{\R}{\mathbb{R}}
\renewcommand{\d}[1]{\,\mathrm{d}#1}
\renewcommand{\div}{\mathrm{div}}
\newcommand{\pw}{\mathrm{pw}}
\newcommand{\s}{\mathrm{s}}
\title{Discrete weak duality of hybrid high-order methods for convex minimization problems}
\author{Ngoc Tien Tran}
\date{Source: arXiv:2308.03223v3 (CC BY 4.0)}
\begin{document}
\maketitle

\noindent\emph{Source and licence.} Discrete weak duality of hybrid high-order methods for convex minimization problems. Version arXiv:2308.03223v3 (CC BY 4.0), distributed by arXiv under the Creative Commons Attribution 4.0 International licence, http://creativecommons.org/licenses/by/4.0/, as recorded in the arXiv licence metadata for this exact version and shown on the abstract page https://arxiv.org/abs/2308.03223v3. Full licence notice: \texttt{licenses/arxiv-2308.03223-CC-BY-4.0.txt}.\par\smallskip

\noindent\textit{Editorial scope.} Counted unit: the article's theorem thm:error-estimate (source0.tex lines 461--469). The article's complete original proof of it is delivered with it (source0.tex lines 470--495, one \texttt{proof} environment of 26 lines). No mathematics is written, repaired or re-derived: the statement and the proof are the article's own lines, in the article's own notation, and no retained line is substituted or re-flowed.  Retained uncounted as the source-local context the proof needs: lines 255--257; lines 268--270; lines 273--275; lines 446--448; lines 457--459.  Nothing was omitted from any retained span, so the package carries no omission record.  No retained line contains a citation, so the wrapper carries no bibliography and no bibliography entry.  No retained line contains a figure environment, an image-inclusion command or drawing code, so no image is delivered and the package declares no asset dependency.  Editorial formatting only, in addition: an AMS \texttt{amsart} preamble that restates the article's own declarations and macros used by the retained lines, namely the environments \texttt{lemma}, \texttt{theorem} on separate counters, together with the standard packages those lines need.  The retained environments print in this document as Lemma 1, Theorem 1 in order of appearance, whereas the article prints the same items with the numbering of its own full document.  The complete native source of the article as distributed in the arXiv e-print for this exact version is bundled unchanged as \texttt{sources/145281/source0.tex}.
\begin{lemma}[commuting property]\label{lem:commuting-property}
	Any $v \in V$ and $\tau \in W^{1,1}(\Omega;\M)$ satisfy $\D_h \I_{V} v = \Pi_\Mcal^k \D v$, $\div_h \I_W \tau = \Pi_\Mcal^k \div \tau$, and the orthogonality $\D_\pw(v - \mathcal{R}_h \I_V v) \perp \D_\pw P_{k+1}(\Mcal;\R^m)$ with respect to the $L^2$ scalar product.
\end{lemma}

\begin{align}\label{def:dual-pot-rec-1}
	\int_\Omega \mathcal{R}_h^* \tau_h : \D_\pw v_{k+1} \d{x} = - \int_\Omega v_{k+1} \cdot \div_h \tau_h \d{x} + \sum_{S \in \Sigma} \int_S [v_{k+1}]_S \cdot \tau_S \d{s}
\end{align}

\begin{align}\label{def:dual-pot-rec-2}
	\Pi_{Z(\Mcal)} \mathcal{R}_h^* \tau_h = \Pi_{Z(\Mcal)} \tau_\Mcal.
\end{align}

\begin{align}\label{def:lower-order-terms}
	\psi(x,a) \coloneqq -f(x) \cdot a \quad\text{and}\quad \psi_h(x,a) \coloneqq -f_h(x)\cdot a
\end{align}

\begin{align}\label{def:strong-duality}
	\max E^*(W_\mathrm{N}) = E^*(\sigma) = E(u) = \min E(V_\mathrm{D}) < \infty.
\end{align}

\begin{theorem}[a~priori]\label{thm:error-estimate}
	Suppose that \eqref{def:lower-order-terms}--\eqref{def:strong-duality} hold. If $\sigma \in W^{1,1}(\Omega; \M)$, then any $\xi \in L^{p'}(\Omega; \M)$ with
	$\xi \in \partial \Psi(\D_h \I_V u)$ and $\varrho \in L^p(\Omega; \M)$ with $\varrho \in \partial \Psi^*(\mathcal{R}_h^* \I_W \sigma)$ a.e.~in $\Omega$ satisfy
	\begin{align}\label{ineq:error-estimate}
		E_h(\I_V u) &- \min E_h(V_\mathrm{D}(\Mcal)) \leq -\int_\Omega (1 - \Pi_\Mcal^k) \xi:(1 - \Pi_\Mcal^k) \D u \d{x}\nonumber\\
		& + \int_\Omega (\D_h \I_V u - \varrho) : (\sigma - \mathcal{R}_h^* \I_W \sigma) \d{x} + \int_\Omega (f - f_h) \cdot (u - \mathcal{R}_h \I_V u) \d{x}\nonumber\\
		& + \s_h(\I_V u)/r + \gamma_h(\I_W \sigma)/r' - \sum_{S \in \Sigma \setminus \Sigma_\mathrm{N}} \int_S [\mathcal{R}_h \I_V u]_S \cdot (1 - \Pi_S^k)\sigma \nu_S \d{x}.
	\end{align}
\end{theorem}

\begin{proof}
	The proof departs from the split
	\begin{align}\label{ineq:proof-error-estimate-split}
		E_h(\I_V u) - E_h(u_h) \leq E_h(\I_V u) - E(u) + E^*(\sigma) - E_h^*(\I_W \sigma).
	\end{align}
	Since $0 \leq \Psi(\D u) - \Psi(\D_h \I_V u) - \xi : (\D u - \D_h \I_V u)$ a.e.~in $\Omega$ from $\xi \in \partial \Psi(\D_h \I_V u)$, the projection property $\D_h \I_V u = \Pi_\Mcal^k \D u$ from \Cref{lem:commuting-property} imply
	\begin{align}\label{ineq:proof-convergence-rates-primal}
		E_h(\I_V u) - E(u) \leq - \int_\Omega (\xi : (1 - \Pi_\Mcal^k) \D u - (f - f_h) \cdot u) \d{x} + \s_h(\I_V u)/r&.
	\end{align}
	The property $\varrho \in \partial \Psi^*(\mathcal{R}_h^* \I_W \sigma)$ provides $0 \leq \Psi^*(\sigma) - \Psi^*(\mathcal{R}^*_h \I_W \sigma) - \varrho : (\sigma - \mathcal{R}^*_h \I_W \sigma)$ a.e.~in $\Omega$. Since $\div_h \I_W \sigma = \Pi_\Mcal^k \div \sigma = -f_h$ from \Cref{lem:commuting-property}, this shows
	\begin{align}\label{ineq:proof-convergence-rates-dual}
		E^*(\sigma) - E^*_h(\I_W \sigma) \leq - \int_\Omega \varrho : (\sigma - \mathcal{R}_h^* \I_W \sigma) \d{x} + \gamma_h(\I_W \sigma)/r'.
	\end{align}
	The $L^2$ orthogonality
	$\D_h \I_V u - \D_\pw \mathcal{R}_h \I_V u \perp \D_\pw P_{k+1}(\Mcal;\M)$ from \Cref{lem:commuting-property}, $\D_h \I_V u \in P_k(\Mcal;\M)$, and $\Pi_{Z(\Mcal)}(\Pi_\Mcal^k \sigma - \mathcal{R}_h^* \I_W \sigma) = 0$ from \eqref{def:dual-pot-rec-2} lead to
	\begin{align}\label{eq:proof-error-estimate-ibp-1}
		\int_\Omega (\sigma - \mathcal{R}_h^* \I_W \sigma) : \D_h \I_V u \d{x} = \int_\Omega (\sigma - \mathcal{R}_h^* \I_W \sigma) : \D_\pw \mathcal{R}_h \I_V u \d{x}.
	\end{align}
	The definition \eqref{def:dual-pot-rec-1} of $\mathcal{R}_h^*$, a piecewise integration by parts, $\div \sigma = - f$, and $\div_h \I_W \sigma = - f_h$ prove
	\begin{align}\label{eq:proof-error-estimate-ibp-2}
		\int_\Omega (\sigma &- \mathcal{R}_h^* \I_W \sigma) : \D_h\I_V u \d{x}\nonumber\\
		&\qquad = \int_\Omega (f - f_h) \cdot \mathcal{R}_h \I_V u \d{x}
		+ \sum_{S \in \Sigma \setminus \Sigma_\mathrm{N}} \int_S [\mathcal{R}_h \I_V u]_S \cdot (\sigma - \Pi_S^k \sigma) \nu_S \d{x}.
	\end{align}
	The combination of this with \eqref{ineq:proof-error-estimate-split}--\eqref{ineq:proof-convergence-rates-dual} concludes the proof.
\end{proof}

\end{document}
